\clearpage
\SourceLeaf{73}%
\pdfbookmark[1]{Göttingische gelehrte Anzeigen. 1808. December 5.}{review1808}
\begin{center}\small Göttingische gelehrte Anzeigen. 1808. December 5.\end{center}
\medskip
Das vor kurzem erschienene Programm, worin Hr. Prof. \textsc{Gauss} seine Vorlesungen
für das gegenwärtige Winter-Halbjahr ankündigt, hat vornehmlich zum
Zweck, die Astronomen auf die Vortheile aufmerksam zu machen, welche man
aus den jetzt mit Unrecht zu sehr vernachlässigten \textit{Sternhöhen} für Zeit- und Breitenbestimmungen
ziehen kann. Es ist zwar nicht zu läugnen, dass diese Gattung
von Beobachtungen, besonders mit Reflexions-Instrumenten, schwieriger
ist, und mehr Übung erfordert, als Sonnen-Beobachtungen: indess ist es eben
so gewiss, dass diese Schwierigkeiten nicht unüberwindlich sind, und dass die
darauf gewandte Mühe sich in sehr vielen Fällen vielfach belohnt. Vorzüglich
ist die Anwendung der Sternhöhen reisenden Beobachtern und Seefahrern sehr
zu empfehlen, nicht bloss desswegen, weil sie die oft vereitelten Sonnenhöhen
oftmals ersetzen können, sondern auch, weil jene einen wichtigen Vortheil geben
können, den man entbehren muss, wenn man sich auf letztere einschränkt.
Durch einzelne Sonnenhöhen kann man nur dann die Zeitbestimmung erhalten,
wenn man die Breite schon kennt, und umgekehrt muss jene als bekannt vorausgesetzt
werden, wenn man die Breite durch eine Sonnenhöhe ausser dem
Meridian bestimmen will. Durch die Verbindung zweier Sonnenhöhen kann
man allerdings Zeit und Breite zugleich nach bekannten Methoden bestimmen:
allein, wenn das Resultat einige Schärfe haben soll, müssen jene Beobachtungen
\SourceLeaf{74}%
nothwendig ziemlich weit von einander entfernt sein, wobei also eine gute Uhr
ein wesentliches Bedürfniss ist, und zur See ausserdem noch die Bewegung des
Schiffes in Betrachtung gezogen werden muss. Dagegen kann man aus der Verbindung
der beobachteten Höhe eines Sternes mit der (wenn man will, unmittelbar
nachher) beobachteten Höhe eines zweiten Sterns allemal Zeit und Breite zugleich
mit grosser Schärfe bestimmen, wenn nur die Sterne selbst schicklich gewählt
sind: der Gang der Uhr braucht dabei nur so weit zuverlässig bekannt zu
sein, dass man die Zwischenzeit zwischen den beiden Beobachtungen in Sternzeit
verwandeln kann. Die Auflösung dieses Problems lässt sich sogleich ohne
Mühe auf die Auflösung dreier sphärischer Dreiecke zurückführen, wobei indess
der numerische Calcul etwas weitläuftig wird. Hr. Prof. \textsc{Gauss} hat aber eine
merklich kürzere Auflösung bloss auf analytischem Wege entwickelt, welche zugleich
als eine Bestätigung angesehen werden kann, dass alles, was aus geometrischen
Betrachtungen geschlossen wird, immer wenigstens eben so einfach und
elegant durch die Analyse gefunden werden kann, wenn diese auf eine schickliche
Art angewandt wird. Auch zum practischen Gebrauch wird man die gegenwärtige
Auflösung weit bequemer finden, als diejenige welche Hr. \textsc{Kraft} im
13. Bande der \textit{Nova Acta Petropol.} für dasselbe Problem (nur mit der, jedoch
unwesentlichen, Einschränkung, dass die beiden Höhenmessungen gleichzeitig
sein sollen) gegeben hat: eine noch bequemere indirecte Auflösung verspricht der
Verfasser bei einer andern Gelegenheit mitzutheilen. Überdiess enthält die gegenwärtige
Abhandlung die Entwickelung der Umstände, auf welche man bei der
Wahl der Sterne zu sehen hat, damit die Resultate möglichst genau ausfallen,
und die Erläuterung des gelehrten Verfahrens durch die wirkliche Anwendung
auf ein paar auf der hiesigen Sternwarte gemessene Sternhöhen, aus denen die
Polhöhe mit der von \textsc{Tob. Mayer} festgesetzten auf $7''$ übereinstimmend, und die
Zeitbestimmung nur um eine halbe Zeitsecunde verschieden von der an demselben
Tage auf andere Art ausgemittelten abgeleitet wird.
\par\bigskip\centerline{\rule{.24\textwidth}{.3pt}}
